% !TEX root = ./../main.tex

\section{Conclusion and Future Work}
\label{sec:conclusion}

The objective of this work is to explore the use of Refinement Types to formally reason about \ac{COP} 
programs, and manage the possible inconsistencies or incompleteness that may arise at run time, in 
response to the interaction between layers. Refinements Types are chosen due to their ease of 
integration with existing programming languages, and their powerful features that allow the programmer 
to formulate functions' pre-conditions and post-conditions, as well as to prove mathematical statements 
about the program's functionality. There are many ways to implement \ac{COP}, however, expressing it 
as a set of partial functions that can be composed together to incrementally build programs, significantly 
simplifies the process of writing and verifying \ac{COP} applications.

To validate the usefulness of Refinement Types in verifying \ac{COP} programs, we use the Bellevue 
graphical editor, as a validating case study implemented in two iterations. The first iteration evidences 
errors caused by unforeseen interactions between the layers in \ac{COP} programs. Modifications to the 
program, such as changing the composition order of predicated functions, or adding a condition to the 
activation predicates, can introduce bugs that are hard to detect, as the decoupling between application 
logic and layer activation makes it harder to reason about the program. The second iterations uses
Refinement Types, with the abstractions to mange them, making the bugs explicit and statically 
verifiable. Enforcing the assumptions stated in the activation predicates into the application's logic, 
inconsistencies and incompleteness can be avoided while defining the functions and their composition. 
Once this is guaranteed, the programmer can then specify more complex properties to verify the 
interactions between functions by describing the dependencies between them in terms of activation and 
state manipulation.

As presented in \fref{sec:discussion}, the \scode{PredicatedFunction} and \scode{Chainable} 
abstractions could be improved, and more work could be done for verifying properties that involve 
states in different time instants. In particular, there are two main possibilities to improve this work.
First, we can make the \scode{andThen} and \scode{orElse} functions depend on a \scode{Chainable} 
argument (or a similar abstraction), in such a way that we disallow composing two functions unless there 
is evidence that the two functions can in fact be composed (in this case, by virtue of the 
\scode{isChainable} property). Currently, if programmers omit a \scode{Chainable} instance to verify a 
layer interaction, they can still compose the layers; the composition functions are unrestricted. Achieving 
this level of abstraction is challenging, because there may be other properties besides 
\scode{isChainable} that we would like to specify as a pre-requirement for composing two functions.
Hence, this new \scode{andThen} function would need to receive arbitrarily user defined predicates, 
instead of constant ones like the \scode{isChainable} predicate.
Second, extending the current abstractions would open the possibility to specify properties about the 
state at different points in time. We discussed in \fref{sec:discussion}, that a good example of this would 
be to statically guarantee conditions of a predicated behavior with respect to specific events 
(\eg activation of a layer) at a previous (not necessarily adjacent) time.  A potentially good approach 
would be to use Temporal Logic with Refinement Types.


\endinput

